\section{Conception du pkjs (téléphone)}
\label{sec:telephone}

\subsection{Composants}

\begin{figure}[H]
  \centering
  \begin{tikzpicture}[
      font=\sffamily\scriptsize,
      c/.style={draw=exigBleu, fill=exigBleu!5, rounded corners=2pt, align=center,
                text width=21mm, minimum height=8mm},
      b/.style={draw=exigGris, fill=exigGris!6, rounded corners=2pt, align=center,
                text width=21mm, minimum height=8mm},
      fleche/.style={-Stealth, draw=exigGris!90!black},
    ]
    \node[b] (msg) at (0, 0)     {J-01 Messagerie};
    \node[c] (vue) at (0, -1.6)  {J-11 Vues};
    \node[c] (hom) at (3.5, 1.2) {J-06 Maison};
    \node[c] (his) at (3.5, 0)   {J-07 Historique};
    \node[c] (ctl) at (3.5, -3.6) {J-10 Commandes};
    \node[c] (cac) at (7.2, 0.6) {J-04 Cache};
    \node[c] (cli) at (7.2, -2.6) {J-03 Client Tado};
    \node[c] (aut) at (10.9, -2.6) {J-02 Liaison};
    \node[b] (ext) at (10.9, -0.6) {Tado};
    \draw[fleche] (msg) -- (vue);
    \draw[fleche] (vue.east) -- ++(0.4,0) |- (hom.west);
    \draw[fleche] (vue.east) -- ++(0.4,0) |- (his.west);
    \draw[fleche] (vue.east) -- ++(0.4,0) |- (ctl.west);
    \draw[fleche] (hom.east) -- ++(0.5,0) |- (cac.west);
    \draw[fleche] (cac) -- (cli);
    \draw[fleche] (cli) -- (aut);
    \draw[fleche] (cli.east) -- ++(0.35,0) |- (ext.west);
  \end{tikzpicture}
  \caption{Composants du pkjs~; une flèche va d'un composant à celui dont il se sert (les services
    J-06, J-07 et J-10 passent tous par le cache et le client)}
  \label{fig:composants-js}
\end{figure}

\begin{table}[H]
  \centering\small
  \begin{tabularx}{\linewidth}{@{}l>{\raggedright\arraybackslash}p{28mm}X>{\raggedright\arraybackslash}p{34mm}@{}}
    \toprule
    \textbf{Réf.} & \textbf{Composant} & \textbf{Responsabilité} & \textbf{Exigences} \\
    \midrule
    J-01 & Messagerie & Reçoit les demandes de la montre, découpe les vues en petits messages, les envoie un à un en attendant l'accusé de chaque envoi. Gère l'attente de l'événement \texttt{ready}. & \refexig{TPB-UTI-005}, \refexig{TPB-ACT-001} \\
    J-02 & Liaison & Conduit le flux de code d'appareil (RFC~8628), conserve les jetons, les rafraîchit, sait les révoquer. Seul composant à connaître les jetons. & \refexig{TPB-CPT-001} à \refexig{TPB-CPT-005}, \refexig{TPB-SEC-001}, \refexig{TPB-SEC-002} \\
    J-03 & Client Tado & Seul composant à connaître les adresses, les formes des requêtes et des réponses. Ajoute le jeton, relit les en-têtes, normalise les erreurs (\S\ref{sec:erreurs}). & \refexig{TPB-ACT-005} \\
    J-04 & Cache & Conserve chaque réponse avec l'heure de son obtention et son TTL, dans le \texttt{localStorage}. Sert les âges. & \refexig{TPB-ACT-002}, \refexig{TPB-INF-001} \\
    J-06 & Maison & Liste des pièces, détail d'une pièce, présence de la maison. & \refexig{TPB-PIE-001} à \refexig{TPB-PIE-010}, \refexig{TPB-DET-001} à \refexig{TPB-DET-007}, \refexig{TPB-PRE-001} à \refexig{TPB-PRE-005} \\
    J-07 & Historique & Rapports du jour, calcul des douze colonnes. & \refexig{TPB-HIS-001} à \refexig{TPB-HIS-007} \\
    J-10 & Commandes & Forçages, retour à la programmation, extinction, présence, raccourcis~; relecture après chaque commande. & \refexig{TPB-PIL-001} à \refexig{TPB-PIL-011}, \refexig{TPB-RAC-001} à \refexig{TPB-RAC-007} \\
    J-11 & Vues & Construit les vues envoyées à la montre~: noms abrégés et distincts, températures en dixièmes dans l'unité du compte, niveaux, indicateurs d'âge, liste des raccourcis offerts. & \refexig{TPB-PLT-006}, \refexig{TPB-AFF-001}, \refexig{TPB-AFF-005} \\
    \bottomrule
  \end{tabularx}
  \caption{Composants du pkjs}
\end{table}

\subsection{Requêtes vers Tado}
\label{sec:requetes}

Le tableau~\ref{tab:requetes} fait le lien entre ce que l'utilisateur fait et les requêtes que cela demande. Les
chemins viennent de [E3] (base \texttt{https://my.tado.com/api/v2}, \texttt{\{h\}} pour l'identifiant de maison,
\texttt{\{z\}} pour celui de pièce)~; ils ont été vérifiés (§\ref{sec:observations}).

\paragraph{Vérification.} Les chemins et schémas du tableau ont été essayés contre le compte de l'utilisateur~: toutes les lectures répondent \texttt{200}, aux formes attendues. Les écarts et les réponses d'écriture sont au §\ref{sec:observations}.

\begin{table}[H]
  \centering\small
  \begin{tabularx}{\linewidth}{@{}>{\raggedright\arraybackslash}p{30mm}>{\raggedright\arraybackslash}p{58mm}>{\raggedright\arraybackslash}p{16mm}X@{}}
    \toprule
    \textbf{Besoin} & \textbf{Requête(s)} & \textbf{TTL} & \textbf{Coût et remarques} \\
    \midrule
    Liste des pièces (F1), détail (F2) & \texttt{GET /homes/\{h\}/zoneStates} & 5 min & 1. Température, puissance, consigne, forçage, fenêtre ouverte, prochain changement, humidité de toutes les pièces. \\
    Présence (F5) & \texttt{GET /homes/\{h\}/state} & 5 min & 1. \texttt{presence} et \texttt{presenceLocked}. \\
    Appareils présents & \texttt{GET /homes/\{h\}/mobileDevices} & 10 min & 1, à l'ouverture de l'écran Présence seulement. Appareils où \texttt{geoTrackingEnabled} et \texttt{atHome}. \\
    Historique (F4) & \texttt{GET /homes/\{h\}/zones/\{z\}/dayReport?date=} & 30 min (jour courant), illimité (jour échu) & 1 par pièce~; 2 par pièce avant midi. \\
    Réglages de la maison & \texttt{GET /homes/\{h\}} & illimité & 1, à la liaison. Unité de température, fuseau (H5). \\
    Plage de consigne & \texttt{GET /homes/\{h\}/zones/\{z\}/capabilities} & illimité & 1 par pièce, à la première ouverture de l'écran de consigne. Observé~: 5 à 25~°C (41 à 77~°F), pas annoncé de 0,1. \\
    Forcer, éteindre (F3) & \texttt{PUT /homes/\{h\}/zones/\{z\}/overlay} & --- & 1, puis relecture \texttt{GET .../zones/\{z\}/state} (1). Total 2. \\
    Retour à la programmation & \texttt{DELETE /homes/\{h\}/zones/\{z\}/overlay} & --- & 1, puis relecture (1). Total 2. \\
    Raccourcis \enquote{tout off} & \texttt{POST /homes/\{h\}/overlay} & --- & 1, puis \texttt{zoneStates} (1). Total 2. \\
    Raccourci \enquote{programmation partout} & \texttt{DELETE /homes/\{h\}/overlay?rooms=} & --- & 1, puis \texttt{zoneStates} (1). Total 2. \\
    Présence imposée & \texttt{PUT /homes/\{h\}/presenceLock} & --- & 1, puis \texttt{state} (1). Total 2. \\
    Présence rendue à Tado & \texttt{DELETE /homes/\{h\}/presenceLock} & --- & 1, puis \texttt{state} (1). Accepté sans Auto-Assist (204)~; \texttt{403} possible pour un autre compte. \\
    Liaison & \texttt{POST login.tado.com/oauth2/device\_authorize}, \texttt{.../token} & --- & Hors de l'API v2~; compté ou non par Tado, ce n'est pas établi. \\
    \bottomrule
  \end{tabularx}
  \caption{Requêtes de l'application, coût et durée de validité}
  \label{tab:requetes}
\end{table}

\subsection{Réponses réelles observées}
\label{sec:observations}

Les requêtes du tableau~\ref{tab:requetes} ont été jouées le 28 octobre 2026 sur le compte de
l'utilisateur (neuf pièces de chauffage, installation classique), les écritures sur une seule pièce
revenue ensuite à sa programmation. Les réponses brutes sont conservées hors du dépôt (elles
contiennent des données personnelles)~; ne figurent ici que les constats qui changent la conception.

\begin{description}
  \item[Quota.] \texttt{RateLimit-Policy} vaut \texttt{"perday";q=20000;w=86400}~; \texttt{RateLimit}
    vaut \texttt{"perday";r=19999} après la première requête. Chaque requête, lecture ou écriture,
    décompte un. Le quota est celui d'un compte Auto-Assist bien que l'utilisateur n'ait pas
    l'abonnement (H2).
  \item[Identifiant de maison.] Il vient de \texttt{GET /me} (une requête de plus que le
    tableau~\ref{tab:requetes}, à faire une fois à la liaison).
  \item[Lectures.] \texttt{zoneStates}, \texttt{state}, \texttt{mobileDevices}, \texttt{zones},
    \texttt{capabilities} et
    \texttt{dayReport} (\texttt{callForHeat.dataIntervals} avec \texttt{NONE}/\texttt{HIGH}…)
    répondent \texttt{200} aux formes décrites.
  \item[Plage de consigne.] \texttt{capabilities} annonce \texttt{celsius} 5 à 25 par pas de 0,1 et
    \texttt{fahrenheit} 41 à 77 par pas de 0,1.
  \item[Présence.] \texttt{PUT presenceLock} avec \texttt{\{"homePresence":"HOME"\}} répond \texttt{204}~;
    \texttt{state} indique alors \texttt{presenceLocked = true} et ajoute
    \texttt{showSwitchToAutoGeofencingButton = true}. \texttt{DELETE presenceLock} répond \texttt{204}
    \emph{sans abonnement Auto-Assist}~; \texttt{presenceLocked} revient à \texttt{false}.
  \item[Forçage.] Le corps du \texttt{PUT overlay} est
    \texttt{\{"setting":\{"type":"HEATING","power":"ON","temperature":\{"celsius":17\}\},
    "termination":\{"typeSkillBasedApp":"TIMER","durationInSeconds":1800\}\}}. Tado répond
    \texttt{200} avec le forçage complet (\texttt{setting}, \texttt{termination} avec
    \texttt{expiry} et \texttt{remainingTimeInSeconds})~: la relecture par \texttt{GET .../state} n'est
    pas indispensable pour connaître le résultat. \texttt{TIMER}, \texttt{TADO\_MODE} et
    \texttt{MANUAL} sont acceptés tels quels~; \texttt{power = OFF} met la température à \texttt{null}.
    \texttt{NEXT\_TIME\_BLOCK} est converti par Tado en \texttt{TIMER} expirant au créneau suivant~;
    \texttt{TADO\_MODE} est donc la valeur à envoyer pour \enquote{prochain créneau}.
  \item[Fahrenheit.] Une consigne \texttt{\{"fahrenheit":62\}} est acceptée directement~; Tado la
    convertit en 16,5~°C (61,7~°F) et la renvoie dans les deux unités. La valeur effective peut donc
    différer de celle demandée~: la montre affiche la valeur relue (réponse du \texttt{PUT}), non celle
    qui a été saisie.
  \item[Retour à la programmation.] \texttt{DELETE overlay} répond \texttt{204}~; la pièce reprend
    sa consigne programmée.
\end{description}

\subsection{Politique de lecture}
\label{sec:lecture}

\begin{enumerate}
  \item Toute lecture passe par le cache~: si la donnée a moins de son TTL, aucune requête.
  \item Après une commande, la relecture ignore le cache (\refexig{TPB-ACT-005}).
  \item L'action \enquote{actualiser} de l'utilisateur (\refexig{TPB-ACT-003}) ignore le cache de la
        donnée de l'écran en cours, et d'elle seule.
  \item L'actualisation automatique pendant la consultation (\refexig{TPB-ACT-004}) vaut un
        \texttt{zoneStates} toutes les cinq minutes tant que l'écran des pièces est affiché.
\end{enumerate}

\subsection{Règles de calcul}

\subsubsection{Historique}
\label{sec:calcul-historique}

Pour chaque pièce, le rapport du jour donne des périodes d'appel de chaleur
(\texttt{callForHeat.dataIntervals}, chacune avec un début, une fin et une valeur parmi
\texttt{NONE}, \texttt{LOW}, \texttt{MEDIUM}, \texttt{HIGH}). L'historique affiche les $n$ dernières heures
pleines (douze, huit sur Time Round), la dernière étant l'heure en cours.

\begin{enumerate}
  \item Les colonnes sont les heures locales de la maison (H5). Si l'heure locale est avant
        midi, le rapport de la veille est aussi demandé.
  \item L'intensité d'une colonne est $\sum_i d_i\,w(v_i) / D$, où $d_i$ est la durée de la période $i$
        recouvrant la colonne, $w$ le poids de DEC-11 (0, 25, 50, 100~\%) et $D$ la durée écoulée de
        la colonne (une heure, ou le temps écoulé pour l'heure en cours).
  \item Elle est rangée dans la plage $[0;10[$, $[10;25[$, $[25;50[$, $[50;100]$ (niveaux 0 à 3). Une valeur à
        la limite va à la plage supérieure (\refexig{TPB-HIS-002}).
  \item Si les périodes du rapport ne couvrent pas au moins 90~\% de la durée écoulée de la
        colonne, la colonne est \emph{inconnue}, non nulle (\refexig{TPB-HIS-007}).
  \item La colonne de l'heure en cours est marquée (\refexig{TPB-HIS-003}).
\end{enumerate}

\subsection{Commandes}

\subsubsection{Forçage d'une pièce}

\begin{table}[H]
  \centering\small
  \begin{tabular}{@{}ll@{}}
    \toprule
    \textbf{Choix de l'utilisateur} & \textbf{Condition de fin envoyée à Tado} \\
    \midrule
    30 min, 1 h, 2 h & \texttt{TIMER}, \texttt{durationInSeconds} = 1800, 3600, 7200 \\
    Prochain créneau & \texttt{TADO\_MODE} \\
    Jusqu'à nouvel ordre & \texttt{MANUAL} \\
    Éteindre le chauffage & \texttt{power = OFF}, \texttt{TADO\_MODE} \\
    \bottomrule
  \end{tabular}
  \caption{Correspondance entre durées et conditions de fin}
\end{table}

La température est envoyée en degrés Celsius. Lorsque l'unité du compte est le degré
Fahrenheit, la montre propose des degrés entiers et le pkjs convertit (Tado accepte aussi les degrés Fahrenheit, voir ci-dessous). Le dernier forçage d'une pièce (\refexig{TPB-INF-002}) est mémorisé par le pkjs
dans son stockage, après confirmation de Tado.

Tado accepte une consigne exprimée en degrés Fahrenheit (§\ref{sec:observations})~: le pkjs peut l'envoyer telle quelle.

\subsubsection{Présence et raccourcis}
\label{sec:presence}

\paragraph{Mode de présence.} Le pkjs déduit le mode de \texttt{state}~:
\texttt{presenceLocked = false} donne \emph{Auto}~; \texttt{true} donne \emph{Présent} ou
\emph{Absent} selon \texttt{presence}. Imposer un mode est \texttt{PUT presenceLock}~; revenir à
\emph{Auto} est \texttt{DELETE presenceLock}. Un refus \texttt{403} de ce dernier signale l'absence
d'abonnement Auto-Assist~: le pkjs l'annonce, mémorise que \emph{Auto} n'est pas
disponible, et ne l'offre plus (H2).

\paragraph{Raccourcis offerts.} Le pkjs construit la liste à chaque ouverture de l'écran.

\begin{table}[H]
  \centering\small
  \begin{tabularx}{\linewidth}{@{}>{\raggedright\arraybackslash}p{40mm}>{\raggedright\arraybackslash}p{50mm}X@{}}
    \toprule
    \textbf{Raccourci} & \textbf{Offert si} & \textbf{Effet} \\
    \midrule
    Tout off & toujours & \texttt{POST /overlay}~: toutes les pièces à \texttt{power = OFF}, \texttt{MANUAL} (jusqu'à nouvel ordre, comme l'application Android). \\
    Activer la programmation dans toutes les pièces & au moins une pièce a un forçage (\texttt{overlay} non nul dans \texttt{zoneStates}) & \texttt{DELETE /overlay?rooms=} avec les pièces forcées. \\
    Mode Absent & \texttt{presence = HOME} et \texttt{presenceLocked = false} (la géolocalisation identifie l'utilisateur comme présent) & \texttt{PUT presenceLock} \texttt{AWAY}. \\
    Mode Géolocalisation & \texttt{presenceLocked = true}, Auto-Assist non exclu & \texttt{DELETE presenceLock}. \\
    \bottomrule
  \end{tabularx}
  \caption{Raccourcis et conditions d'offre}
\end{table}

\subsection{Liaison au compte Tado}
\label{sec:liaison}

\begin{enumerate}
  \item La montre demande la liaison (au premier lancement, ou après une liaison perdue).
  \item J-02 appelle \texttt{device\_authorize} avec \texttt{scope=offline\_access} et reçoit
        \texttt{user\_code}, \texttt{verification\_uri\_complete}, \texttt{expires\_in} (300~s) et
        \texttt{interval} (5~s).
  \item J-01 envoie à la montre le code et son délai~; la montre l'affiche (écran Liaison). Le
        pkjs ouvre aussi sur le téléphone la page de validation.
  \item J-02 interroge \texttt{token} toutes les cinq secondes (valeur de \texttt{interval}) jusqu'à
        recevoir les jetons, un refus, ou l'expiration.
  \item Les jetons sont écrits dans le \texttt{localStorage}, puis J-06 lit les réglages de la maison
        (H5) et le choix de la maison est fait (\refexig{TPB-CPT-002}).
  \item À chaque usage, J-02 vérifie l'âge du jeton d'accès~; s'il est périmé, il rafraîchit,
        \emph{écrit le nouveau jeton de rafraîchissement avant de poursuivre}, puis rejoue la requête
        une seule fois.
\end{enumerate}

Une requête qui reçoit un refus d'accès après rafraîchissement déclenche la liaison perdue
(\refexig{TPB-CPT-005}). Délier (\refexig{TPB-CPT-004}) efface les jetons, le cache et le
stockage de la montre.
